\subsection{Effective Saturating Potential}
  \label{subsec:effective_potential}

  The field equation introduced below is mathematically identical to
  the quasi-linear MOND formulation in the deep galactic regime.
  Here, this structure is not postulated as a fundamental modification of gravity.
  Rather, we argue that the empirical MOND phenomenology emerges
  as the effective low-density limit of a more general saturation
  principle governing kinematic inference.

  We define the effective gravitational potential $\Phi_{\mathrm{eff}}$ through a
  modified Poisson equation of the form
  \begin{equation}
    \nabla \cdot \left[
                   \mu\!\left( \frac{|\nabla \Phi_{\mathrm{eff}}|}{a_\star} \right)
                   \nabla \Phi_{\mathrm{eff}}
    \right]
    =
    4 \pi G \rho_{\mathrm{b}} ,
    \label{eq:modified_poisson}
  \end{equation}
  where $\rho_{\mathrm{b}}$ denotes the baryonic mass density
  and $\mu(x)$ encodes the saturation of the effective response.

  For clarity, we work primarily with the $\mu$-formulation in
  Eq.~\eqref{eq:modified_poisson}.
  An equivalent algebraic form is sometimes written using an interpolation
  function $\nu(y)$ defined by $g_{\mathrm{eff}} = \nu(y) g_{\mathrm{N}}$
  with $y = g_{\mathrm{N}}/a_\star$.
  In that notation, $\nu(y) \rightarrow 1$ for $y \gg 1$,
  ensuring an exactly Newtonian limit.

  In the present interpretation, $\mu(x)$ does not represent a fundamental
  field-dependent coupling.
  It parametrizes the progressive saturation of the effective response
  as the underlying transition structure approaches its maximal resolvable rate.
  The specific functional form adopted here provides a minimal interpolation
  between regimes and is not intended as a fundamental law.

  The function $\mu(x)$ satisfies the limiting behaviors
  \begin{equation}
    \mu(x) \rightarrow 1 \quad \text{for} \quad x \gg 1 ,
  \end{equation}
  and
  \begin{equation}
    \mu(x) \rightarrow x \quad \text{for} \quad x \ll 1 .
  \end{equation}

  We assume $\mu(x)$ to be positive, monotone increasing, and at least $C^{1}$,
  so that $\Phi_{\mathrm{eff}}$ and its first derivatives remain continuous
  across the transition.
  This regularity avoids spurious discontinuities in the effective force
  and ensures stable orbital dynamics.

  In the high-acceleration regime,
  Eq.~\eqref{eq:modified_poisson} reduces to the standard Poisson equation.
  In the low-acceleration or diffuse regime,
  the effective response saturates,
  leading to a finite asymptotic acceleration.

  For spherically symmetric systems,
  the effective radial acceleration $g_{\mathrm{eff}}(r)$ satisfies
  \begin{equation}
    \mu\!\left( \frac{g_{\mathrm{eff}}}{a_\star} \right)
    g_{\mathrm{eff}}
    =
    g_{\mathrm{N}} ,
    \label{eq:effective_acceleration}
  \end{equation}
  where $g_{\mathrm{N}}$ is the Newtonian acceleration sourced by baryons alone.

  Equation~\eqref{eq:effective_acceleration} yields asymptotically flat rotation
  curves for isolated galaxies with finite baryonic mass.
  The transition scale $a_\star$ controls both the onset of flattening
  and the amplitude of the asymptotic velocity.

  Importantly, $a_\star$ is treated as a universal parameter.
  No halo-by-halo fitting or galaxy-dependent adjustment is introduced.
  All diversity in rotation curve shapes arises from the observed baryonic
  distributions.
  In this work, $a_\star$ is interpreted as a minimal resolvable acceleration
  within the effective description,
  reflecting an intrinsic bound on the transition rate
  of the geometric response.
